\subsection*{一份恢复台账对应一份可继续的责任}

现在把具体过程接起来。设本地恢复基点 $b_0=99$，J 在 100 登记 E1，101 保存 $A_J$，102 保存 $R_J$，103 关闭 E1，104 冻结 g7。恢复者已经取得完整 1 至 104 的权威前缀，其中快照覆盖至 99。平台保留 X1，履约库保留 F501。

恢复者在第 9 代建立台账，记录 g7、$b_0$、$b_{\rm cut}=104$、恢复来源版本和物化游标。它按次序处理 100 至 104，把 E1 的固定截止与关闭状态一起恢复，按 101 找回 P1/C1/T1，按 102 恢复完整原结果、O1 和 EV1，再把会话最终状态定为关闭。重放不能在处理到 100 的 \code{active} 状态时提前开放，否则 103 的不可撤回关闭尚未生效。

\CFourteenFigure{ch14-original12}{从 99 恢复到 104 的过程中，活动登记只是中间投影，第 103 位关闭应用以后才取得最终旧会话状态}{fig:ch14-prefix-replay}

台账对每个 C 保存原受理位置、平台范围和键、证据状态、外部路径状态、本地物化状态以及下一责任。C1 可以标为成功已物化。另一个 C3 若返回 \code{accepted\_pending/H3}，就标为已封闭首次受理、沿 H3 待终结。C4 若只有查空而没有终结依据，就标为外部未决、禁止自动新发送，并保存负责恢复者和下次查询安排。

游标只跨过已经逐项处置的连续决定。处置可以是完成物化，也可以是把未决 C 及其固定关联完整转入可扫描台账。跳过一个损坏的受理记录再把游标记成 104，会让系统声称已经覆盖却无法枚举这份责任。遇到无法解释的记录，恢复保持冻结或只开放已经独立隔离、证据完整的其他分区。

\subsection*{旧重试在三种入口各得到什么}

第一种是购买提交入口。E1 已关闭，返回 410 和 \code{operation\_expired}，说明旧准入结束。若能够经授权定位 P1，可以给出其查询位置。恢复尚未建好 P1 时，响应提供账户授权的恢复查询入口，购买编号待原受理关联恢复以后取得。过期响应只说明这次旧准入已经结束，原购买的资金状态从查询和恢复入口取得。

第二种是已受理购买的查询入口。它根据授权后的完整 S、原受理回执或 P1 读取 J 与本地投影，组合可解释的一份结果。J 已有原成功，而本地投影尚未补齐时，可以返回原受理及 \code{restoring} 状态，随后提供原 \code{payment\_result}，也可以直接从可信结果来源返回固定结果。接口要把当前 \code{progress} 与原成交结果分开。

第三种是服务端 C/T 恢复入口。它从已经受理的 C1 取得原范围和 \code{pay-C1}，依据平台当前保护条件选择查询、原身份安全重试或原子关闭。E1 的关闭结束旧会话准入，C1 的未决责任保存在 J 和恢复台账。读取不到本地 C1 时，先从 J 重建原关联，再执行外部动作。

\CFourteenFigure{ch14-entry-branches}{同一旧购买在提交、授权查询和服务端恢复入口分别检查资格、原结果与外部路径}{fig:ch14-entry-branches}

下面是本例在 E1 已关闭、C1 原成功已经恢复后的查询结果。\code{acceptance} 固定原受理，\code{progress} 反映查询时的业务进度，\code{payment\_result} 保留原付款决定：

\begin{BookCode}
{
  "result_version": "purchase-result-external-v2",
  "acceptance": {
    "purchase_id": "P1",
    "charge_id": "C1",
    "operation": {
      "account_id": 1,
      "action": "purchase.create.external.v2",
      "session_id": "E1",
      "idempotency_key": "K1"
    },
    "product_id": "paper-001",
    "quantity": 1,
    "quote_id": "Q7",
    "currency": "CNY",
    "amount_minor": 3000
  },
  "progress": {
    "business_state": "paid",
    "fulfillment_state": "fulfillment_accepted"
  },
  "payment_result": {
    "purchase_id": "P1",
    "charge_id": "C1",
    "order_id": "O1",
    "status": "succeeded",
    "provider_scope": "simulator:merchant-7:test:charge.v1",
    "provider_charge_id": "X1",
    "currency": "CNY",
    "amount_minor": 3000
  }
}
\end{BookCode}

\code{fulfillment\_accepted} 来自履约库保留的 F501 及可信回报经过购买库受控入口恢复后的状态。只恢复 $R_J$、尚未核对这份回报时，O1 的履约进度为 \code{unfulfilled} 或接口公布的恢复中状态，资金成功字段保持。付款成功、建 F 与实际发货继续分别取证。

具有恢复资格的同一 S 异参请求在已恢复绑定处返回 \code{conflict}。查到 J 中完整前缀却没有该 S，处理也要结合会话资格：E1 已关闭便拒绝旧输入，新会话有效且归属开放，才允许首次受理。若 J 覆盖不完整，空结果保留恢复中的含义。同一个本地空查询，根据权威覆盖与准入状态落入不同分支。

\subsection*{什么时候可以恢复受理}

恢复开放至少包含四类具体判据。首先，g7 的承载者和代次已经由唯一权威决定，旧地域无法在规则之外创建新受理，资源端已确认相应写入水位。其次，恢复者已经取得完整受理集合、原输入、会话资格和不会重用的编号来源，所有旧 S 都能被识别或被关闭资格拦住。

再次，每项旧 C 都有持续保存的责任：成功已经固定并可物化，最终拒绝具有覆盖证据，未决操作则保存原平台身份和下一步，所有可能迟到或未来成功的路径都受平台原身份保护、不可延长截止或同域关闭约束。最后，本地 P/C/O/结果/EV 的归属关系已校验，消费者的 I/B/F 及 \code{dispatch} 从自己的保存域恢复，跨域缺口有明确的冻结范围。

这些条件容许一部分支付继续等待。C3 可以沿 H3 长期等待，只要首次受理已经封闭、结果通道可继续查询、责任保存在台账中，新地域就沿原付款身份恢复它。此时可以受限开放独立新意图，并对未决总量设置容量限制。若 C4 的旧到达可能越过已经结束的去重窗口再次执行，相关支付方式的自动发送和准入保持受控，直到取得能够覆盖该路径的依据。

重新开放的实际检查还包括台账与数据库的原子边界。物化游标和对应行共同提交，重启后才能从已知边界继续。开放模式的改变在权威域提交，入口每次首次受理检查当前模式。仅在管理页面点亮绿色状态，而入口继续使用旧缓存，会把上述检查留在业务执行以外。

\CFourteenFigure{ch14-original13}{受限开放依据唯一归属、完整身份与资格、受保护的旧外部路径，以及可恢复的本地和消费关系}{fig:ch14-reopen}

\subsection*{逐个迁移怎样保持原承诺}

最后把本章的安全关系写在同一个状态上。设 J 保存权威受理与结果决定，$\Gamma$ 保存会话资格和登记进度，D 保存归属路由代次与模式（本章的 D 指目录，扣款流水在资金模型中另按其作用域解释），W 是购买库投影，L 是平台账本，Q 是履约保存域。对每个完整 S，J 至多保存一份不可覆盖的原输入与 P/C 绑定。同 C 的所有外部尝试使用原范围和原平台身份，同一资金成功固定唯一 O 和 EV，初次履约继续在原 B 上恢复实际 F。

空初态没有受理和外部效果，这些关系成立。$A_J$ 只在唯一权威仲裁中建立一次 S 绑定，前置业务拒绝保存没有付款关联的稳定结果，重试和异参保持原记录。E 受平台原身份规则约束，同 C 的可能到达汇到同一平台操作。成功 $R_J$ 共同保存可信资金事实、固定 O/EV 和原结果，身份最终拒绝则保存终结依据与原拒绝结果。重复证据恢复已有决定。本地 A/R 重放只复制已有决定，并以事务约束核对全套关系。消费者按原 I/B 取得 F501，重复 EV 恢复原初次授权。

冻结和路由代次授予调整哪个地域能够推进工作，权威受理日志中的原身份及平台键随原购买保存。会话 E1 关闭后，旧提交入口结束准入，已经接受的支付和任务则由日志 J 与恢复台账接续。购买库投影 W 回退时，入口进入 \code{recovering}，恢复者从继续保存的受理、会话与目录历史重建完整图像。后继入口取得相应覆盖依据以后，再对新的首次受理作出判断。每一项允许迁移都维持单一绑定、原结果和对应效果身份的关系。

这份归纳的外部前提是平台持续保护所有可能生效的原身份到达，或提供足够的封闭与终结证据。若 J 和所有独立关联也随旧备份丢失，恢复者会重新遇到图~\ref{fig:ch14-indistinguishable}~的两条不可区分历史。若旧平台请求可以在身份遗忘以后再次执行，外部效果上限也会失去前提。对应入口通过冻结和持久核查限制后续动作。

进展则来自一组可运行的路径。权威域最终能提交目录与 J 的决定，恢复来源能够提供完整前缀，物化数据库能够提交，恢复者持续获得执行机会。成功的 X 最终能被查询或对账发现，平台已受理工作能终结并保留可取回结果，消费者和回报通道最终恢复处理。客户端能够持久保留原身份并继续查询，才会最终获知自己的原结果。网络持续分区、证据永久丢失或外部操作永久未决时，相应责任保持在台账中。

地域切换完成以后，Alice 的旧购买继续是 E1/K1 对应的 P1/C1，支付继续是原范围内的 \code{pay-C1/X1}，订单继续是 O1，事件继续是 EV1，初次履约继续指向 501。Alice 之后确实再买一份，才为新的业务意图登记 \code{E9:1}、持久保存 K2，并取得另一份购买与平台身份。最终系统的说明便可以据这些边界，写出各种故障以后的恢复入口、暂停受理条件和未决工作的持续责任。
\FloatBarrier
